\section{Path-valued inversion of the original raw density}
\label{sec:v52-path-inversion}
A path limit under a fixed roof window does not by itself control the
conditional path law at the individual roof values inside that window.
We estimate the path numerator in a norm which permits the path test to
be chosen after the roof value. The estimate retains the complete
physical remainder. In particular, the unresolved positive height is
not replaced by a local integral.

Let $\mathcal C=C([0,1],\R^4)$ with its supremum metric, and let
$\mathcal F$ consist of the real functions $F$ on $\mathcal C$ satisfying
$\|F\|_\infty\le1$ and $\operatorname{Lip}(F)\le1$. For a finite signed
measure $\sigma$ on $\mathcal C$ put
\[
 \|\sigma\|_{\mathrm{BL}^*}=\sup_{F\in\mathcal F}|\sigma(F)|.
\]
On probability measures this is $d_{\rm BL}$. On positive finite
measures it is exactly the mass. This is not total variation in the
path variable. Write $\mathsf W_R=\operatorname{Law}(\mathbb B_{\Omega_R})$
for the collision bridge law in Section~\ref{sec:v41-bridges}.

The original exact-label source is
$\one_{\Pi_{n,k,m,R}}\nu_R^*$, and its roof density is $p_{n,R}(k,m,u)$.
Disintegrate this finite source with respect to its roof. Let
$\mathsf Q_{m,R}^{n,k,u}$ be the resulting conditional law of
$\mathcal B_{m,R}$, the collision path tied to its own terminal record
in \eqref{eq:v41-collision-bridge}. For almost every roof value define
finite path measures
\begin{align}
 \mathbf P_{m,R}^{n,k}(u)&=m^2p_{n,R}(k,m,u)\mathsf Q_{m,R}^{n,k,u},
                                      \label{eq:v52-path-density}\\
 \mathbf b_{B,m,R}^{n,k}(u)(F)
  &=m^2b^{\varepsilon(B),F(\mathcal B_{m,R})}_{n,k,m,R}(u),
 \qquad \varepsilon(B)=A_0B^{-1/12}.    \label{eq:v52-bad-path-density}
\end{align}
The second measure is obtained by disintegrating the positive source
$(1-H_{m,R}^{\varepsilon(B)})\one_{\Pi_{n,k,m,R}}\nu_R^*$ jointly
with the first one. Thus
$0\le\mathbf b_B(u)\le\mathbf P(u)$ as measures, and its mass is
$\mathfrak b_B(u)=m^2b^{\varepsilon(B),1}_{n,k,m,R}(u)$.
We take both measures to be zero when their scalar density is zero.
The source factor $1/c$ is already in $p$ and $b$ and is not inserted
again. The scalar reference is $G_{m,R}^{n,k}(u)$ of
Section~\ref{sec:v51-positive-remainder}.

These disintegrations exist on the standard Borel source and path
spaces. Constancy of discrete prefixes in the endpoint argument below
is only a statement on one protected word chart, not across a physical
or section boundary. A countable class $\mathcal F_0\subset\mathcal F$ norms
$\mathrm{BL}^*$ on all finite measures on $\mathcal C$: choose a
countable dense set of paths, approximate the admissible values on
its finite subsets by rationals, and extend them with the same
Lipschitz bound, followed by truncation to $[-1,1]$. Diagonal
approximation gives pointwise convergence to every member of
$\mathcal F$, and bounded convergence gives the norming assertion.
Consequently all roof norms below are measurable, and all estimates
can be imposed off one roof-null set for each fixed set of labels.

\subsection{A uniform endpoint budget for the entire path test class}
\begin{lemma}[Lipschitz path tests are protected-admissible]
\label{lem:v52-path-endpoint-budget}
On every regular endpoint chart with the $\varepsilon/2$-margins,
\begin{equation}\label{eq:v52-path-endpoint-budget}
 \operatorname{Lip}_{\rm endpoints}(\mathcal B_{m,R})
                    \le C\varepsilon^{-1}m^{-1/2}.
\end{equation}
Here the image carries the supremum metric of $\mathcal C$.
In particular, uniformly over $F\in\mathcal F$, the insertion
$F(\mathcal B_{m,R})$ has supremum at most one and weak endpoint
gradient at most $C\varepsilon^{-1}m^{-1/2}$.
For each fixed margin, the whole class is eventually admissible
with the common budget $(M,L)=(1,1)$ of
Theorem~\ref{thm:v46-protected-correction}.
\end{lemma}
\begin{proof}
Write the endpoint arclengths as $a,b$ and the internal contact
arclengths as $y_j(a,b)$. The first variation of the prefix roof is
\[
 dS_{j,R}=-p_0\,da+p_j\,dy_j\quad(1\le j<m),\qquad
 dS_{m,R}=-p_0\,da+p_m\,db.
\]
The internal reflection equations cancel all other contact terms in
each prefix. Equations~\eqref{eq:v45-endpoint-influence} and
\eqref{eq:v45-summed-influence}, in their regular version supplied by
Lemma~\ref{lem:v46-regular-endpoints}, give
$|dy_j|\le C\varepsilon^{-1}(|da|+|db|)$ uniformly in $j,m$.
Since $|p_j|\le1$, every prefix roof has endpoint derivative at most
$C\varepsilon^{-1}$. On the same chart every displacement prefix and
every occupation prefix is constant: the physical word and every
section decision persist. Only the roof coordinate of
$\mathsf S_j-(j/m)\mathsf S_m$ can therefore have a nonzero derivative.
After division by $\sqrt m$ its bound is
$C\varepsilon^{-1}m^{-1/2}$. Linear interpolation takes convex
combinations of consecutive grid values and does not enlarge this
supremum bound. Integrate on smaller convex chart squares to obtain
\eqref{eq:v52-path-endpoint-budget}. Composition with a member of
$\mathcal F$ is a scalar Lipschitz function on a two-dimensional
chart; its almost-everywhere gradient has the stated bound.

This argument differentiates neither a measurable path selector nor a
function across a physical or section boundary. For a reconstruction
band $B$ the margin $\varepsilon(B)$ remains fixed before $m$ grows.
The threshold for the common $(1,1)$ budget is permitted to depend on
that fixed margin.
\end{proof}

\subsection{Fixed-band inversion in the path norm}
For $F\in\mathcal F$ let $p^F$ denote the exact-label roof density
with insertion $F(\mathcal B_{m,R})$. The scalar density is $p=p^1$.

\begin{lemma}[Uniform bounded-Lipschitz factorization at fixed band]
\label{lem:v52-fixed-band-path}
For every fixed reconstruction band $B$,
\begin{equation}\label{eq:v52-fixed-band-path}
 \sup_{R,n,k,u,F\in\mathcal F}
 \left|m^2(K_B*p^F)(u)
      -\mathsf W_R(F)G_{m,R}^{n,k}(u)\right|\longrightarrow0.
\end{equation}
The band is fixed in this limit. No estimate for its spectral
constants at $B=B_m$ is included.
\end{lemma}
\begin{proof}
We first explain why the cylinder factorization extends to the whole
class $\mathcal F$, rather than merely to one fixed cylinder.
The Schwartz function $|K_B|$ has an integrable nonnegative
band-limited upper function. For example put
\[
 q(s)=\left(\frac{\sin s}{s}\right)^2+
       \left(\frac{\sin(s+\pi/2)}{s+\pi/2}\right)^2,
\]
with the removable values filled in. The two numerators never vanish
simultaneously, and $q(s)\ge c_0/(1+s^2)$: this follows directly for
$|s|\ge\pi$ and by positivity and compactness on the remaining
interval. With $\widehat f(\xi)=\int_{\R}e^{-i\xi s}f(s)\,ds$, one has
\[
 \widehat q(\xi)=\pi(1-|\xi|/2)_+
                       (1+e^{i\pi\xi/2})\in L^1(\R).
\]
Its support is $[-2,2]$, and inversion uses the factor $(2\pi)^{-1}$.
Thus $\mathsf q^{[B]}(s)=C_Bq(s)\ge\max\{|K_B(s)|,|K_B(-s)|\}$
for a finite constant depending on the fixed reconstruction band.
The bracketed superscript denotes multiplication by $C_B$, not a
rescaling of the Fourier support. This dependence is not used at a
growing band.

At a target $(R,n,k,u)$ consider the signed finite measure on paths
\[
 \sigma_m(F)=m^2\int_{\Pi_{n,k,m,R}}
           F(\mathcal B_{m,R})K_B(u-T_{n,R})\,d\nu_R^*.
\]
Its variation is dominated by the positive path measure obtained
with $\mathsf q^{[B]}(T_{n,R}-u)$. The latter has uniformly bounded mass by the
fixed-band scalar inversion. The derivative calculation in
Proposition~\ref{prop:v41-pinned-fourth}, now with the fixed
nonnegative roof function $\mathsf q^{[B]}$, gives for its grid increments
\begin{equation}\label{eq:v52-band-path-fourth}
 m^2\int_{\Pi_{n,k,m,R}}\mathsf q^{[B]}(T_{n,R}-u)
  |\mathcal B_{m,R}(s)-\mathcal B_{m,R}(r)|^4\,d\nu_R^*
                                   \le C_B|s-r|^2.
\end{equation}
Indeed its Fourier transform is integrable and compactly supported,
so the same fourth derivative of the chronological three-block
product is integrated. Centering at each moving peak cancels the
linear drift, the size factors are $O(l^2)$ for a block of $l$
collisions, and the four-frequency inverse costs $m^{-2}$.
Division by $m^2$ from the fourth power of the path normalization,
followed by the outside factor $m^2$, gives $C_B(l/m)^2$.
The interpolation argument of that proposition gives
\eqref{eq:v52-band-path-fourth} at all real times. The paths start
at zero. The fourth-moment continuity criterion therefore gives
uniform tightness of these positive finite path measures; equivalently,
add mass at the zero path and divide by a common mass upper bound
before applying the probability version. This also proves tightness
of the variations of $\sigma_m$.

For a fixed grid of observation times and a fixed characteristic
cylinder, Lemma~\ref{lem:v41-pinned-factorization}, in its
band-limited proof with $K_B$, gives
\[
 \sigma_m(\phi)-\sigma_m(1)\mathsf W_R(\phi)\longrightarrow0
\]
uniformly in all targets. The exact identities
$\sum_i m_iv_i^{(m)}=0$ cancel both the drift and the mixed Gaussian
term on every moving resonance chart. The exponentially damped
curvature-continuity estimate then replaces the chart Hessian by
$\Omega_R$. All occupation resonances are retained.
Polygonal observation times rather than their grid approximations
change the test by $O(m^{-1/2})$ uniformly, since the one-collision
record is bounded; the dominating mass bound pays that error.

To prove uniformity over $\mathcal F$, suppose it failed and choose
a violating sequence of targets and counts. Pass to a subsequence
with $R\to R_0$ and $\sigma_m(1)\to a$. Lemma~\ref{lem:v53-signed-compactness}, using the variation-tightness
just proved, gives a weakly convergent subsequence of the signed
finite measures. One can obtain it by applying compactness to the
two positive measures $\tau_m+\sigma_m$ and $\tau_m-\sigma_m$,
where $\tau_m$ is their positive dominating measure and
$|\sigma_m|\le\tau_m$ explicitly.
Cylinder characteristic functions at rational times determine finite
signed measures on continuous path space. The factorization identifies
the limit as $a\mathsf W_{R_0}$. Continuity of $\Omega_R$ gives
continuity of $\mathsf W_R$. On a compact set carrying all but an
arbitrarily small amount of the variations, the restrictions of
$\mathcal F$ have a finite uniform net by equicontinuity and boundedness.
It follows that
$\|\sigma_m-\sigma_m(1)\mathsf W_R\|_{\mathrm{BL}^*}\to0$,
contradicting the chosen sequence. Finally the scalar fixed-band
reconstruction gives
$\sup|\sigma_m(1)-G_{m,R}^{n,k}(u)|\to0$. This proves
\eqref{eq:v52-fixed-band-path}.
\end{proof}

\subsection{The path-valued physical remainder}
\begin{theorem}[A common positive remainder for all path tests]
\label{thm:v52-path-remainder}
For every sufficiently large fixed $B$,
\begin{equation}\label{eq:v52-path-remainder}
 \limsup_{m\to\infty}\sup_{R,n,k}\operatorname*{ess\,sup}_{u}
 \left\|\mathbf P_{m,R}^{n,k}(u)
       -G_{m,R}^{n,k}(u)\mathsf W_R
       -\mathbf b_{B,m,R}^{n,k}(u)\right\|_{\mathrm{BL}^*}
                                          \le CB^{-1/192}.
\end{equation}
The same physical incidence and clearance source is used for every
test $F$; it is not selected anew from the test or the roof value.
Consequently, for every fixed $h>0$,
\begin{equation}\label{eq:v52-path-local-variation}
 \sup_{R,n,k,t}\frac1h\int_t^{t+h}
 \left\|\mathbf P_{m,R}^{n,k}(u)
       -G_{m,R}^{n,k}(u)\mathsf W_R\right\|_{\mathrm{BL}^*}\,du
                                               \longrightarrow0.
\end{equation}
Both statements concern the complete original exact-label source.
\end{theorem}
\begin{proof}
For a test $F\in\mathcal F$ use the exact source identity
$p^F=f^{\varepsilon,F}+d^{\varepsilon,\varepsilon,F}
+b^{\varepsilon,F}$. Write $a^F=f^{\varepsilon,F}
+d^{\varepsilon,\varepsilon,F}$. Then exactly
\begin{align*}
 m^2p^F-G\mathsf W_R(F)-m^2b^{\varepsilon,F}
  =m^2(a^F-K_B*a^F)
   +\{m^2K_B*p^F-G\mathsf W_R(F)\}
   -m^2K_B*b^{\varepsilon,F}.
\end{align*}
Fix $B$ and hence $\varepsilon=\varepsilon(B)$.
Lemma~\ref{lem:v52-path-endpoint-budget} makes every insertion in
this test class eventually protected-admissible with budget $(1,1)$.
The protected correction is therefore $O(B^{-1/2})$ in the normalized
essential norm. The complete decision correction is
$O(\varepsilon^{1/16})$ in the same norm by bounded-source domination.
The middle term tends to zero uniformly over the full test class by
Lemma~\ref{lem:v52-fixed-band-path}. The last term is at most
$C\varepsilon^{1/16}\|K_1\|_1$ after the collision limsup, by
Lemma~\ref{lem:v51-smoothed-physical-height}, uniformly over these
bounded insertions. Thus the limsup of the displayed error is at most
$CB^{-1/192}$, uniformly over $F$.

Take the countable norming class $\mathcal F_0$ to make the exceptional
roof set common to all tests. This converts the uniform scalar
estimate into the essential norm in \eqref{eq:v52-path-remainder}.
No derivative of a roof-dependent choice of test is taken.
The positive measure $\mathbf b_B(u)$ has norm equal to its mass.
Hence integration over a translated interval and
\eqref{eq:v49-local-limsup} give a limsup at most
$CB^{-1/192}+C\varepsilon(B)^{1/16}$ in
\eqref{eq:v52-path-local-variation}. Let $B$ tend to infinity only
after this collision limsup. This proves the second assertion.
The local window and the auxiliary spectral band used for its
physical estimate have precisely the order of limits in
Theorem~\ref{thm:v49-full-local-correction}.
\end{proof}

\begin{corollary}[Roof-dependent Lipschitz path selectors]
\label{cor:v52-roof-path-selectors}
The local bound \eqref{eq:v52-path-local-variation} controls, uniformly,
all jointly Borel functions $\Psi(u,x)$ supported in a translated fixed
roof interval for which $\Psi(u,\cdot)\in\mathcal F$ for every $u$.
The selector may depend on $m,R,n,k$ and may be chosen after $u$.
There is no regularity condition in $u$. This assertion is not a
Gaussian-amplitude theorem for arbitrary measurable path selectors.
\end{corollary}
\begin{proof}
At almost every $u$, evaluation against $\Psi(u,\cdot)$ is bounded
by the path dual norm in \eqref{eq:v52-path-local-variation}.
Integrate that bound. The countable norming construction, or bounded
measurable approximation on the standard Borel product, gives the
same conclusion for a measurable selector. The proof first bounds
the whole test class at each roof; it does not differentiate
$\Psi(T_{n,R},\mathcal B_{m,R})$ as a source insertion.
\end{proof}
